% ===== PRÉAMBULE CANONIQUE — à coller VERBATIM en tête de CHAQUE .tex =====
\documentclass[11pt,a4paper]{article}
\usepackage[utf8]{inputenc}\usepackage[T1]{fontenc}\usepackage{lmodern}
\usepackage[french]{babel}
\usepackage{amsmath,amssymb,mathtools}\usepackage{icomma}
\usepackage{siunitx}\sisetup{locale=FR}  % \qty{}{}, \si{}, \ang{}
\usepackage{xcolor}\usepackage{colortbl}
\usepackage{tikz}\usepackage{pgfplots}\pgfplotsset{compat=1.18}
\usepackage[siunitx,europeanresistors]{circuitikz} % circuits (élec.)
\usepackage[most]{tcolorbox}\usepackage{enumitem}\usepackage{array,booktabs}
\usepackage[a4paper,margin=2.2cm]{geometry}
\usetikzlibrary{arrows.meta,calc,angles,quotes,positioning,patterns,%
  backgrounds,shapes.geometric,decorations.pathmorphing,decorations.markings,intersections}
\definecolor{primary}{RGB}{222,49,99}\definecolor{primarydark}{RGB}{150,22,55}
\definecolor{defcol}{RGB}{0,114,178}\definecolor{methcol}{RGB}{52,92,148}
\definecolor{excol}{RGB}{0,135,90}\definecolor{attncol}{RGB}{200,40,55}
\tcbset{boxbase/.style={enhanced,breakable,boxrule=0.9pt,arc=2.5pt,
  left=3mm,right=3mm,top=2.3mm,bottom=2.3mm,fonttitle=\bfseries,coltitle=white}}
% --- boîtes TUTORIEL (noms EXACTS, ne pas renommer) ---
\newtcolorbox{cle}[1][]{boxbase,colback=defcol!6,colframe=defcol,
  title={\textsc{À retenir}\ifx\relax#1\relax\relax\else\textmd{ : \itshape #1}\fi}}
\newtcolorbox{methode}[1][]{boxbase,colback=methcol!7,colframe=methcol,sharp corners,
  title={\textsc{Méthode}\ifx\relax#1\relax\relax\else\textmd{ --- \itshape #1}\fi}}
\newtcolorbox{exemplebox}[1][]{boxbase,colback=excol!5,colframe=excol,
  title={\textsc{Exemple}\ifx\relax#1\relax\relax\else\textmd{ : \itshape #1}\fi}}
\newtcolorbox{piege}[1][]{boxbase,colback=attncol!6,colframe=attncol,
  title={\textsc{Piège}\ifx\relax#1\relax\relax\else\textmd{ : \itshape #1}\fi}}
\newtcolorbox{remarque}[1][]{boxbase,colback=black!4,colframe=black!45,
  title={\textsc{Remarque}\ifx\relax#1\relax\relax\else\textmd{ : \itshape #1}\fi}}
% --- boîtes EXERCICE / CORRIGÉ (noms EXACTS, auto-comptées) ---
\newtcolorbox[auto counter]{exo}[1][]{boxbase,colback=primary!4,colframe=primary,
  title={\textsc{Exercice}~\thetcbcounter\ifx\relax#1\relax\relax\else\textmd{ --- \itshape #1}\fi}}
\newtcolorbox[auto counter]{corrige}[1][]{boxbase,colback=excol!4,colframe=excol,
  title={\textsc{Corrigé}~\thetcbcounter\ifx\relax#1\relax\relax\else\textmd{ --- \itshape #1}\fi}}
\newcommand{\vv}[1]{\vec{#1}}\newcommand{\ds}{\displaystyle}
\newcommand{\R}{\mathbb{R}}\newcommand{\N}{\mathbb{N}}\newcommand{\Z}{\mathbb{Z}}
\begin{document}
\sisetup{per-mode=symbol}

\begin{corrige}[objet placé en $2F$]
\begin{enumerate}[label=\alph*)]
  \item Le rayon par $O$ et le rayon parallèle (qui passe par $F'$) se croisent en \textbf{$2F'$}, de
        l'autre côté de la lentille : l'image se forme en $2F'$.
  \item Image \textbf{réelle}, \textbf{renversée} et de \textbf{même taille} que l'objet
        ($\gamma=-1$). C'est le cas symétrique $2F\leftrightarrow 2F'$.
\end{enumerate}
\end{corrige}

\begin{corrige}[quel rayon émergent est correct ?]
Un rayon qui passe par le \textbf{foyer objet $F$} ressort \textbf{parallèle à l'axe}. Le seul tracé
parallèle à l'axe est le \textbf{rayon 1}. (2 correspond à la règle du rayon parallèle incident ;
3 = non dévié ; 4 est quelconque.)
\end{corrige}

\begin{corrige}[compléter le tableau des natures]
Toutes ces images (sauf le dernier cas) sont \textbf{réelles et renversées} ; seule la taille change.
\begin{enumerate}[label=\alph*)]
  \item image \textbf{réelle}, ponctuelle, au foyer $F'$ ($v=f$).
  \item image \textbf{réelle}, \textbf{renversée}, \textbf{réduite} (entre $F'$ et $2F'$).
  \item image \textbf{réelle}, \textbf{renversée}, \textbf{agrandie} (au-delà de $2F'$) --- le projecteur.
  \item image \textbf{virtuelle}, \textbf{droite}, \textbf{agrandie}, du même côté que l'objet --- la loupe.
\end{enumerate}
\end{corrige}

\begin{corrige}[construction avec une divergente]
\begin{enumerate}[label=\alph*)]
  \item Les rayons émergents \emph{divergent} ; seuls leurs prolongements se croisent \textbf{du même
        côté que l'objet}, entre l'objet et la lentille. L'image s'y forme.
  \item Image \textbf{virtuelle}, \textbf{droite} et \textbf{réduite} : c'est le comportement
        systématique d'une divergente pour un objet réel.
\end{enumerate}
\end{corrige}

\begin{corrige}[l'objet s'approche de la lentille]
\begin{enumerate}[label=\alph*)]
  \item L'image reste \textbf{réelle et renversée}. Elle part de $F'$ (objet à l'infini), s'éloigne
        vers $2F'$ puis file vers l'infini quand l'objet atteint $F$. Sa taille \emph{croît}
        continûment (réduite, puis même taille en $2F$, puis agrandie).
  \item L'image « bascule » : elle devient \textbf{virtuelle, droite et agrandie}, du même côté que
        l'objet (régime loupe). L'objet et l'image se déplacent alors dans le même sens.
\end{enumerate}
\end{corrige}

\end{document}
